\section{Elementary octave: $3{:}3{:}6$ and $3{:}4{:}6$}

The elementary octave uses the atomic-number coordinate as a registry-facing readout. The active fixture is not a nuclear-reaction model. It is a retained-capacity ladder whose rows can be compared with the 118-element name/symbol registry.

\subsection{Declared ladders}

The two elementary rules are:
\[
\FFL_{336}: 3+3\to 6,
\qquad
\FFL_{346}: 3+4\to 6+1_{\mathrm{shed}}.
\]
For a target registry coordinate $Z$, the retained six-capacity depth is
\[
D_6(Z)=\left\lceil \frac{Z}{6}\right\rceil,
\qquad
C_6(Z)=6D_6(Z),
\qquad
\epsilon_Z=C_6(Z)-Z.
\]
The $3{:}3{:}6$ rule has zero surplus per rung. The $3{:}4{:}6$ rule records one sheddic surplus unit per rung. Both rows preserve $\epsilon_Z$ instead of forcing closure.

\subsection{Registry fixture}

The file \texttt{manual-2/data/elementary/element\_registry\_118.csv} contains one row for each atomic-number coordinate from 1 through 118. The first and last rows are:
\begin{center}
\begin{tabular}{@{}rlll@{}}
\toprule
$Z$ & Symbol & Name & Status \\
\midrule
1 & H & Hydrogen & frozen 118-name/symbol fixture \\
118 & Og & Oganesson & frozen 118-name/symbol fixture \\
\bottomrule
\end{tabular}
\end{center}

\aodliteraturenote{The element-name/symbol registry lane is treated as an external authority lane. The fixture cites the IUPAC periodic-table page and keeps atomic-number comparison separate from the core ladder arithmetic.}

\subsection{Scale-comparison rows}

The elementary data lane provides:
\begin{itemize}[leftmargin=2em]
  \item \texttt{fusion\_ladder\_336.csv};
  \item \texttt{fusion\_ladder\_346.csv};
  \item \texttt{stellar\_scaled\_comparison.csv};
  \item \texttt{pubchem\_element\_map.csv}.
\end{itemize}
The stellar-scale file uses labels such as light-nuclei boundary marker, alpha-chain comparison marker, stellar-burning comparison marker, and heavy-element comparison marker. These are comparison labels, not claims that the AOD row is a stellar nucleosynthesis engine.

\subsection{Sample rows}

\begin{center}
\begin{tabularx}{\linewidth}{@{}lrrrrX@{}}
\toprule
Rule & $Z$ & Depth & Capacity & Residual & Status \\
\midrule
$3{:}3{:}6$ & 6 & 1 & 6 & 0 & closed \\
$3{:}3{:}6$ & 8 & 2 & 12 & 4 & boundary residual \\
$3{:}4{:}6$ & 26 & 5 & 30 & 4 & boundary residual; surplus 5 \\
$3{:}4{:}6$ & 118 & 20 & 120 & 2 & boundary residual; surplus 20 \\
\bottomrule
\end{tabularx}
\end{center}
